\documentclass[10pt,a4paper]{amsart}
\usepackage[margin=19mm]{geometry}
\usepackage{amsmath,amssymb,mathtools}
\usepackage[scr=rsfs,cal=euler]{mathalfa}
\usepackage{xfrac}
\usepackage[unicode,hidelinks,bookmarks=false]{hyperref}
\newcommand{\ie}{i.e.\ }
\newcommand{\cf}{cf.\ }
\newcommand{\resp}{respectively\ }
\newcommand{\Subsection}{\subsection}
\providecommand{\nobreakdash}{\nobreak\hbox{-}}
\setlength{\emergencystretch}{2em}
\multlinegap0pt
\usepackage{amssymb}
\usepackage[normalem]{ulem}
\usepackage{xcolor}
\usepackage{enumerate}
\newtheorem{theorem}{Theorem}
\newtheorem{lemma}{Lemma}
\title{The Skew Commutators of Toeplitz and Hankel operators in vector-valued Hardy Space}
\author{Priyanka Aroda$^*$ and Santanu Dey$^{**}$}
\date{Source: arXiv:2609.04752v1 (CC BY 4.0)}
\begin{document}
\maketitle

\noindent\emph{Source and licence.} The Skew Commutators of Toeplitz and Hankel operators in vector-valued Hardy Space. Version arXiv:2609.04752v1 (CC BY 4.0), distributed by arXiv under the Creative Commons Attribution 4.0 International licence, http://creativecommons.org/licenses/by/4.0/, as recorded in the arXiv licence metadata for this exact version and shown on the abstract page https://arxiv.org/abs/2609.04752v1. Full licence notice: \texttt{licenses/arxiv-2609.04752-CC-BY-4.0.txt}.\par\smallskip

\noindent\textit{Editorial scope.} Counted unit: the article's theorem S+S* (source0.tex lines 1006--1046). The article's complete original proof of it is delivered with it (source0.tex lines 1047--1236, one \texttt{proof} environment of 190 lines). No mathematics is written, repaired or re-derived: the statement and the proof are the article's own lines, in the article's own notation, and no retained line is substituted or re-flowed.  Retained uncounted as the source-local context the proof needs: lines 756--765; lines 997--1005.  Nothing was omitted from any retained span, so the package carries no omission record.  No retained line contains a citation, so the wrapper carries no bibliography and no bibliography entry.  No retained line contains a figure environment, an image-inclusion command or drawing code, so no image is delivered and the package declares no asset dependency.  Editorial formatting only, in addition: an AMS \texttt{amsart} preamble that restates the article's own declarations and macros used by the retained lines, namely the environments \texttt{theorem}, \texttt{lemma} on separate counters, together with the standard packages those lines need.  The retained environments print in this document as Theorem 1, Lemma 1 in order of appearance, whereas the article prints the same items with the numbering of its own full document.  The complete native source of the article as distributed in the arXiv e-print for this exact version is bundled unchanged as \texttt{sources/209386/source0.tex}.
        \begin{theorem}\label{p-3}
           A bounded linear operator $X$ on $H^2$ satisfies 
           $$SX=X^*S$$
           if and only if 
           $$X=U^*\begin{pmatrix}
                 T^*_{\phi_{22}}&T_{\overline{f}+zf}\\
                 0&T_{\phi_{22}}
             \end{pmatrix}U,$$
           for some $ \phi_{22}\in H^{\infty}, f \in H^2 \text{ with } {\overline{f}+zf} \in L^\infty$ and $U$ is defined as above.
        \end{theorem}

        \begin{lemma}\label{special}
            Let $X \in \mathcal{B}(H^2).$ Then 
            $$XS^2={S^*}^2X$$ if and only if $$X=U^*
                \begin{pmatrix}
                H_{\phi_{11}}&H_{\phi_{12}}\\
                H_{\phi_{21}}&H_{\phi_{22}}
                \end{pmatrix}U,$$ 
                for some $\phi_{11}, \phi_{12}, \phi_{21}, \phi_{22} \in L^{\infty}.$
        \end{lemma}

        \begin{theorem}\label{skew is S+S*}
           A bounded operator $Y\in \mathcal{B}(H^2\oplus H^2)$ 
           satisfies
           \begin{equation*}
               (S\oplus S^*)Y=Y^*(S\oplus S^*)
           \end{equation*}
           if and only if 
           \begin{equation}\label{converse skew}
               Y=
               \begin{pmatrix}
                   U^*T_{11}U& U^*\hat{T}_{12}U\\
                   U^*{\hat{T}_{21}}U& U^*T_{22}U
               \end{pmatrix}
           \end{equation}
           where 
           $$T_{11}=
             \begin{pmatrix}
                 T^*_{\phi_{11}}&T_{\overline{f}+zf}\\
                 0&T_{\phi_{11}}
             \end{pmatrix},
             T_{22}=
             \begin{pmatrix}
                T^*_{\phi_{22}}&T_{\overline{g}+zg}\\
                 0&T_{\phi_{22}}
             \end{pmatrix},
             $$
                for some $\phi_{11}, \phi_{22} \in H^\infty, f, g \in H^2$ with ${\overline{f}+zf}, {\overline{g}+zg} \in L^\infty,$ and
             $$
                \hat{T}_{12}=
                \begin{pmatrix}
                0&H_{\hat{\phi}_{12}}\\
                0&0
                \end{pmatrix},
                \hat{T}_{21}=
                \begin{pmatrix}
                0&H^*_{\hat{\phi}_{12}}\\
                0&0
                \end{pmatrix},
            $$
             for some ${\hat{\phi}_{12}} \in L^\infty.$
        \end{theorem}

        \begin{proof}
           Let us write $Y$ in the form of an operator matrix in the decomposition $H^2\oplus H^2$.
           \begin{equation*}
            Y=
               \begin{pmatrix}
                   Y_{11}&Y_{12}\\
                   Y_{21}&Y_{22}
               \end{pmatrix}.
           \end{equation*}
           The above operator equation leads to
           \begin{equation*}
                \begin{pmatrix}
                   SY_{11}&SY_{12}\\
                   S^*Y_{21}&S^*Y_{22}
               \end{pmatrix}
               =
               \begin{pmatrix}
                   Y^*_{11}S&Y^*_{21}S^*\\
                   Y^*_{12}S&Y^*_{22}S^*\\
               \end{pmatrix}.
           \end{equation*}
             By Theorem \ref{p-3}, 
             \begin{align*}
                 Y_{11}=U^*T_{11}U, \text{ where } T_{11}=
             \begin{pmatrix}
                 T^*_{\phi_{11}}&T_{\overline{f}+zf}\\
                 0&T_{\phi_{11}}
             \end{pmatrix},
             \end{align*}
             for some $\phi_{11} \in H^\infty, f \in H^2$ with $\overline{f}+zf \in L^\infty.$
             \begin{align*}
                 Y_{22}=U^*T_{22}U, \text{ where } T_{22}=
             \begin{pmatrix}
                 T^*_{\phi_{22}}&T_{\overline{g}+zg}\\
                 0&T_{\phi_{22}}
             \end{pmatrix},
             \end{align*}
             for some $\phi_{22} \in H^\infty, g \in H^2$ with $\overline{g}+zg\in L^\infty.$
             
             It remains to find the solutions to the following operator equations:
             \begin{equation}\label{equation-1}
                 SY_{12}=Y^*_{21}S^*,
             \end{equation}
             and 
             \begin{equation}\label{equation-2}
                 S^*Y_{21}=Y^*_{12}S.
             \end{equation}
            Composing equation (\ref{equation-1}) with $S$ we get
            \begin{equation*}
                SY_{12}S=Y^*_{21}.
            \end{equation*}
            Equivalently
            \begin{equation}\label{equation-3}
                Y_{21}=S^*Y^*_{12}S^*.
            \end{equation}
            Substituting in equation (\ref{equation-3}) in equation \ref{equation-2} we get
            \begin{equation} \label{equation-4}
                {S^*}^2Y^*_{12}S^*=Y^*_{12}S.
            \end{equation}
            Composing equation (\ref{equation-4}) with $S$, we get
            \begin{equation*}
                 Y^*_{12}S^2={S^*}^2Y^*_{12}.
            \end{equation*}
            Equivalently
            \begin{equation*}
                 Y_{12}S^2={S^*}^2Y_{12}.
            \end{equation*}
            By Lemma \ref{special}, we get
            $Y_{12}=U^*\hat{T}_{12}U,$ where
            \begin{equation*}
            \hat{T}_{12}=
                \begin{pmatrix}
                H_{\hat{\phi}_{11}}&H_{\hat{\phi}_{12}}\\
                H_{\hat{\phi}_{21}}&H_{\hat{\phi}_{22}}
                \end{pmatrix}, 
            \end{equation*}
             for some $\hat{\phi}_{11}, \hat{\phi}_{12}, \hat{\phi}_{21}, \hat{\phi}_{22} \in L^{\infty}.$
            Let $\hat T_{21}=UY_{21}U^*.$ Using the same analogy, we get
            \begin{align*}
                \hat T_{21}=UY_{21}U^*&= U{S}^*Y^*_{12}S^*U^*\\
                &=\tilde{S}^*{\hat{T}_{12}}^*\tilde{S}^*\\
                &=\begin{pmatrix}
                0&I\\
                S^*&0
            \end{pmatrix}
            \begin{pmatrix}
                H^*_{\hat{\phi}_{11}}&H^*_{\hat{\phi}_{21}}\\
                H^*_{\hat{\phi}_{12}}&H^*_{\hat{\phi}_{22}}
                \end{pmatrix}
            \begin{pmatrix}
                0&I\\
                S^*&0
            \end{pmatrix}\\
            &=\begin{pmatrix}
                H^*_{\hat{\phi}_{22}}S^*&H^*_{\hat{\phi}_{12}}\\
                S^*H^*_{\hat{\phi}_{21}}S^*&S^*H^*_{\hat{\phi}_{11}}
            \end{pmatrix}.
            \end{align*}
            We note that the equations (\ref{equation-1}) and (\ref{equation-2}) is equivalent to the following equations:
            \begin{align}\label{converted equation 1}
                \tilde S \hat T_{12}=\hat T^*_{21}{\tilde S}^*,
            \end{align}
            and
            \begin{align}\label{converted equation 2}
                 {\tilde S}^* \hat T_{21}=\hat T^*_{12}\tilde S.
            \end{align}
            Note that equation (\ref{converted equation 1}) implies
            \begin{align*}
                \begin{pmatrix}
                0&S\\
                I&0
            \end{pmatrix}
            \begin{pmatrix}
                H_{\hat{\phi}_{11}}&H_{\hat{\phi}_{12}}\\
                H_{\hat{\phi}_{21}}&H_{\hat{\phi}_{22}}
                \end{pmatrix}
                =\begin{pmatrix}
                H^*_{\hat{\phi}_{22}}S^*&H^*_{\hat{\phi}_{12}}\\
                S^*H^*_{\hat{\phi}_{21}}S^*&S^*H_{\hat{\phi}_{11}}
                \end{pmatrix}^*
                \begin{pmatrix}
                    0&I\\
                    S^*&0
                \end{pmatrix}.
            \end{align*}
            This implies
            \begin{align*}
               \begin{pmatrix}
                SH_{\hat{\phi}_{21}}&SH_{\hat{\phi}_{22}}\\
                H_{\hat{\phi}_{11}}&H_{\hat{\phi}_{12}}
                \end{pmatrix}=
                \begin{pmatrix}
                SH_{\hat{\phi}_{21}}SS^*&SH_{\hat{\phi}_{22}}\\
                H_{\hat{\phi}_{11}}SS^*&H_{\hat{\phi}_{12}}
            \end{pmatrix}.
            \end{align*}
            Therefore we have
            \begin{align}\label{E1}
                SH_{\hat{\phi}_{21}}=SH_{\hat{\phi}_{21}}SS^*, ~H_{\hat{\phi}_{11}}=H_{\hat{\phi}_{11}}SS^*.
            \end{align}
            Hence $H_{\hat{\phi}_{11}}=H_{\hat{\phi}_{21}}=0.$
            \begin{align}\label{E2}
                SH_{\hat{\phi}_{22}}=SH^*_{\hat{\phi}_{22}},
            \end{align}
            By equation (\ref{converted equation 2}) we have
            \begin{align*}
                \begin{pmatrix}
                    0&I\\
                    S^*&0
                \end{pmatrix}
                \begin{pmatrix}
                H^*_{\hat{\phi}_{22}}S^*&H^*_{\hat{\phi}_{12}}\\
                0&0
                \end{pmatrix}
                =\begin{pmatrix}
                0&0\\
                 H^*_{\hat{\phi}_{12}} &H^*_{\hat{\phi}_{22}}
                \end{pmatrix}
                \begin{pmatrix}
                    0&S\\
                    I&0
                \end{pmatrix}.
            \end{align*}
            Thus
            \begin{align*}
                \begin{pmatrix}
                    0&0\\
                    S^*H^*_{\hat{\phi}_{22}}S^*&S^*H^*_{\hat{\phi}_{11}}
                \end{pmatrix}=
                \begin{pmatrix}
                    0&0\\
                    H^*_{\hat{\phi}_{22}}&H^*_{\hat{\phi}_{12}}S
                \end{pmatrix}.
            \end{align*}
            Therefore we get $H_{\hat{\phi}_{22}}=0.$
            Thus
            \begin{align*}
                \hat{T}_{12}=
                \begin{pmatrix}
                0&H_{\hat{\phi}_{12}}\\
                0&0
                \end{pmatrix},
                \hat{T}_{21}=
                \begin{pmatrix}
                0&H^*_{\hat{\phi}_{12}}\\
                0&0
                \end{pmatrix},
            \end{align*}
            for some ${\hat{\phi}_{12}} \in L^\infty.$ Conversely let $Y$ be of the form as defined in equation (\ref{converse skew}). The diagonal equations follows from Theorem \ref{p-3}, and the off diagonal parts are immediate by definition of Hankel operator.
       \end{proof}

\end{document}
